%==============================================================================
%  CHAPTER 3 -- METHODOLOGY   (OPTIONAL)
%
%  Include this chapter ONLY if you developed something yourself: a new
%  algorithm, a new processing chain, a new identification or calibration
%  procedure.
%
%  This is where your ORIGINAL CONTRIBUTION lives. Everything here must be
%  yours; anything taken from the literature belongs in Chapter 2.
%
%  If the work is purely experimental -- you applied an existing method to a
%  new case -- delete this file and its \include line in Main.tex. LaTeX
%  renumbers the remaining chapters automatically.
%
%  Make it REPRODUCIBLE: a competent reader should be able to reimplement the
%  method from this chapter alone.
%==============================================================================
\chapter{Methodology}
\label{ch:methodology}

%------------------------------------------------------------------------------
\section{Rationale}
\label{sec:method-rationale}
%  Why existing methods were not sufficient. Refer back to the limitations you
%  established in Section~\ref{sec:state-of-the-art}.

%------------------------------------------------------------------------------
\section{Proposed method}
\label{sec:method-proposed}
%  The method itself, step by step -- a numbered list, a block diagram or
%  pseudo-code. State every assumption explicitly, and say when it holds.

%------------------------------------------------------------------------------
\section{Implementation}
\label{sec:method-implementation}
%  Software, libraries, parameter choices and the reasoning behind them.
%  If the code is public, give the repository here.

%------------------------------------------------------------------------------
\section{Validation strategy}
\label{sec:method-validation}
%  How you will show the method works: synthetic data with a known answer,
%  comparison against a reference technique, or both.
